\scanpage{34}
\mspage{15}

Mais un élément inversible ne peut être minimal. En effet, on aurait pour tt $u\in M$ un $u'\in M$ tel que $\varphi=uu'$, \deleted{i.e.\ $\unsure{u}$} \unsure{en|donc|et en} posant $v=u'\varphi^{-1}$, on aurait
\[
uv=1
\]
montrons que l'on aurait aussi
\[
p=vu=1,
\]
en effet on a $p^2=p$, i.e. $p.p=\mathrm{id}.p$ donc $p=\mathrm{id}$. Donc tous les objets de $M$ seraient inversibles, i.e. $M$ serait un groupe, ce qui n'est pas le cas.

\uline{NB} Considérons \deleted{une partie} \inserted{(sous les conditions du corollaire) un} sous-$M$-ensemble à droite $F$ de $M$, donc $F$ \deleted{s'identifie}\inserted{peut se} regarder comme un objet de $Y^\wedge$ ($Y=B_M$), d'où une catégorie $X=Y_{/F}$. On peut montrer que les $X$ sont tous $1$-connexes.

J'ai eu du mal à construire un $F$ \deleted{tel que}\inserted{$0$-connexe} qui ne soit $1$-connexe !

% archive sheet number: 33
